\section{REVIEW METHODOLOGY}
\label{sec:review_methodology}

\subsection{Review Design, Protocol, and Reporting Framework}

This work was conducted as a PRISMA 2020-grounded narrative systematic review
with a scoping-style Population--Concept--Context component
\cite{PRISMA2020}. Search reporting was organized in accordance with the
traceability principles of PRISMA-S \cite{PRISMAS2021}. An internally versioned
protocol, decision log, progress tracker, and documented protocol amendments
were maintained throughout the review. The review was not externally
registered.

The review addressed optical and photonic systems in which sensing and
communication functions were jointly considered, integrated, co-designed,
co-optimized, or evaluated within the same architecture, optical link,
waveform or resource framework, hardware platform, channel model, or
application scenario. The primary synthesis frame covered fiber,
free-space-optical, VLC/LiFi, photonic-terahertz, and hybrid optical systems.

\subsection{Information Sources and Search Strategy}

The search window extended from 1 January 2020 through 22 June 2026. Final
searches were executed on 22 June 2026 in Scopus, IEEE Xplore,
ScienceDirect, SpringerLink, Wiley Online Library, and Taylor \& Francis
Online. The search combined integrated-sensing-and-communication terminology
with optical-medium and platform terms, including optical and photonic
systems, fiber or fibre, free-space optical, optical wireless, VLC/LiFi, and
photonic-terahertz variants. Platform-specific refinements were used where
required by search-interface limits.

Exact executed queries, field restrictions, document and language limits,
execution dates, exported record counts, and source filenames were retained in
the versioned search log. Date eligibility was determined from the earliest
verifiable publisher online or public availability date. Records first made
available after 22 June 2026 were excluded even when a database year filter
retrieved the complete 2026 publication year.

\subsection{Eligibility Criteria and Corpus Roles}

The primary technical evidence corpus comprised English-language,
peer-reviewed journal articles, identifiable early-access journal articles,
and full-length conference or proceedings papers published within the search
window. To be eligible, a report had to provide sufficient technical
information to establish an optical or photonic sensing--communication
relationship and to support at least one of the review's synthesis domains,
such as architecture, integration mechanism, metric reporting, validation,
trade-off evidence, or application mapping.

Reports were excluded from the primary technical corpus when they addressed
RF-only ISAC without an optical component, communication without a sensing
function, sensing without communication-system integration, or no genuine
sensing--communication integration. Abstract-only items, posters without a
full technical paper, editorials, opinions, and reports without sufficient
English technical content for eligibility assessment were also excluded.
Review articles and other records retained for background, terminology,
taxonomy cross-checking, or technology lineage were assigned to a contextual
corpus and were not counted as primary technical evidence.

\subsection{Selection Process and Counting Units}

Bibliographic records, reports, and studies were maintained as distinct
counting units. Deduplication preceded title-and-abstract screening. Reports
advanced from screening were sought in full text and assessed against the
eligibility criteria. Full-text exclusions were assigned one primary reason
code with a source-grounded note. Contextual-only reports were retained as a
separate disposition rather than being counted as either full-text exclusions
or included primary evidence.

Screening, full-text assessment, report-to-study mapping, extraction, and
survey-use adjudication were maintained in structured audit ledgers. The
workflow was investigator-supervised and AI-assisted, and investigator-provided
or investigator-delegated adjudication was retained in the provenance record.
Independent duplicate human screening, independent full-corpus human verification, and
third-reviewer arbitration were not documented and are therefore not claimed.

After report-level eligibility was locked, related reports were mapped to the
same underlying study when source evidence supported report equivalence or a
companion-version relationship. Each included study cluster was assigned one
accessible primary extraction report, while companion reports were retained
for bibliographic and evidence lineage. This procedure prevented multiple
reports of one study from being counted as multiple studies.

\subsection{Data Extraction}

Extraction was organized at study, report, and claim levels. Study-level
fields covered bibliographic identity, optical modality, system architecture,
integration mechanism, shared hardware, waveform or resources, sensing task,
communication function, validation type, scenario, application domain,
enabling technologies, limitations, and availability information. Quantitative
records retained the reported metric name, metric role, value, unit,
measurement plane, operating condition, validation context, and source
location. Trade-off records were retained only when both sides of the claimed
relationship and an auditable source location were present.

Reported, calculated, and digitized values were distinguished in the data
model. No project-level graph digitization or newly calculated performance
value was introduced. Missing values were retained as missing rather than
inferred from adjacent text, figures, or related reports.

Claim-level comparison and conflict rules are stated in the companion
claim-governance fragment. In brief, quantitative comparison required
compatible metric definitions, metric roles, measurement planes, scenarios,
and validation conditions. Source conflicts were preserved and governed rather
than silently averaged or resolved without evidence.

% CLAIM-GOVERNANCE INTEGRATION MARKER:
% Insert 03_CLAIM_GOVERNANCE_EN.tex here if a single consolidated Methods
% section is preferred. The paragraph above is already complete and does not
% imply missing Phase-D content.

\subsection{Technical Appraisal and Certainty Assessment}

All 206 included studies underwent a deterministic Technical Quality
Assessment Framework (TQAF) appraisal. Eight dimensions---technical relevance,
metric clarity, reporting completeness, validation maturity, reproducibility,
benchmark readiness, comparison admissibility, and limitation
transparency---were scored on an ordinal 0--3 scale. A ninth field summarized
overall evidence contribution. Each final score was accompanied by a
dimension-level audit record containing the source fields, rule, pre-cap score,
applicable cap, final score, and rationale.

Prespecified guardrails prevented strong aggregate ratings from masking weak
core evidence. An overall score of 3 required technical relevance, metric
clarity, reporting completeness, and validation maturity all to be at least 2.
Studies with both metric clarity and validation maturity at or below 1 were
capped at an overall score of 1. Any study contributing a quarantined claim was
capped at an overall score of 2; studies contributing a quarantined metric
claim were additionally capped at 1 for comparison admissibility. Absence of
an author-reported limitation capped limitation transparency at 1, while a
material source conflict capped it at 2. These rules qualified the language
and weight of synthesis conclusions but did not exclude an otherwise eligible
study.

Certainty was summarized for 115 prespecified S1--S7 evidence bodies by
combining contributing-study count, optical-modality coverage, the median of
the core appraisal dimensions, and the share of studies with adequate core
scores. Broad mixed or unclassified fallback bodies were retained for coverage
and audit only, marked non-substantive, and assigned unclear certainty. The
assessment was deterministic and auditable; it was not an independently
duplicated human quality assessment.

\subsection{Synthesis Strategy}

The included literature was heterogeneous across optical modality, system
architecture, metric definition, measurement plane, scenario, and validation
method. A conventional effect-size meta-analysis was therefore not performed.
The synthesis was structured, descriptive, taxonomy-based, and
metric-governed. Quantitative summaries were admissible only within explicitly
compatible evidence strata, and claim eligibility for quantitative use does
not by itself establish cross-study comparability or meta-analytic
poolability.

Seven synthesis domains were evaluated: optical modality (S1), architecture
and integration mechanisms (S2), metric reporting (S3), communication--sensing
trade-offs (S4), validation maturity and benchmark readiness (S5), enabling
technologies and applications (S6), and evidence gaps and 6G relevance (S7).
Modality, validation maturity, data availability, code/model availability, and
6G relevance were assigned once per study through a frozen 206-row
normalization crosswalk. Integration mechanisms, validation types, enabling
technologies, and applications were multi-label variables; their category
totals were therefore not summed as mutually exclusive prevalence estimates.
The ``other'' fallback was assigned only when no recognized category existed
for the relevant study and axis, while unmatched source tokens were retained
in a normalization audit.

The primary synthesis used 8,203 claims: 3,020 qualitative evidence claims,
4,779 quantitative metric claims, and 404 trade-off claims. The 31 context-only
metric claims and 72 quarantined claims were excluded from primary numerical
synthesis. Claim counts and unique-study counts were reported separately, and
the denominator appropriate to each analysis was stated explicitly.
